% TMLR submission skeleton — F2 methodology paper
%
% Replace [...] placeholders with content from papers/f2_methodology.md.
% The body math is already in dollar-delimited LaTeX in the source;
% prose, tables, and figure references need template-specific markup.
%
% TMLR template source: https://www.overleaf.com/latex/templates/tmlr-journal-submission/fcjxjtfqrsqb

\documentclass{article}

% Use the `tmlr` package: \usepackage{tmlr}    for under-review submissions
%                       \usepackage[accepted]{tmlr}   for camera-ready
\usepackage{tmlr}

\usepackage{amsmath, amssymb, amsthm}
\usepackage{booktabs}
\usepackage{graphicx}
\usepackage{hyperref}
\usepackage{url}

\title{F2: Stratified Pearl-CDE Mediation Analysis for Transformer
Training-Recipe Ablations}

% Double-blind: keep anonymous until accepted
\author{Anonymous Authors\\
\texttt{anonymous@openreview.net}}

\begin{document}
\maketitle

\begin{abstract}
% Paste from papers/f2_methodology.md "Abstract" section.
% Target 250 words. Headline finding: RmsNorm NDE flips sign across
% strata; framework is open-source, reproducible.
[Abstract goes here]
\end{abstract}

\section{Introduction}\label{sec:intro}
% §1.1 The problem with seed-mean ablation
% §1.2 Contributions
% §1.3 Roadmap
[Introduction goes here]

\section{Background}\label{sec:background}
% §2.1 Pearl causal mediation
% §2.2 Sensitivity analysis
% §2.3 ML ablation practice
[Background goes here]

\section{The F2 framework}\label{sec:framework}
% §3.1 Stratification mechanism — Stratum enum, ModeKind registry, Figure 2
% §3.2 Zhao-Luo four-path decomposition — TE = NDE + NIE_M1 + NIE_M2 + NIE_chain
% §3.3 Bridge-score sensitivity envelope — (Γ, Λ), tipping-point classification
% §3.4 Cross-stratum comparator — stable_across_strata flag
% §3.5 Provenance and reproducibility — W3C-PROV preamble + schema-bump discipline
[Framework section goes here]

\section{Sandbox ablation matrix}\label{sec:sandbox}
% §4.1 Setup — Table 3 sandbox configuration
% §4.2 Why sandbox-scale
[Sandbox section goes here]

\section{Empirical results}\label{sec:results}
% §5.1 Canonical ablation — suppression pattern
% §5.2 wd0 stratum — sign flip (headline)
% §5.3 Cross-stratum stability — rms invariant
% §5.4 Sensitivity envelope
[Results section goes here]

% Figure 1 — the headline
\begin{figure}[t]
  \centering
  \includegraphics[width=0.7\linewidth]{figures/fig1_rms_nde_signflip.png}
  \caption{RmsNorm NDE across three strata. The canonical NDE
    (\$-4.12\$ BPB, helpful) reverses sign in the wd0 Pearl-CDE
    stratum (\$+0.43\$ BPB, harmful). Both CIs exclude zero.}
  \label{fig:signflip}
\end{figure}

\section{Sensitivity to design choices}\label{sec:sensitivity}
% §6.1 Mediator pair (M_1, M_2) — Loop 50 swap
% §6.2 Statistic family — t vs perm vs BCa vs bootstrap-t
% §6.3 Strata — choice criteria
[Sensitivity section goes here]

\section{Limitations}\label{sec:limitations}
% 6 limitations: sandbox scale, N=5, no interaction, synthetic task,
% two-mediator only, no post-treatment confounders
[Limitations section goes here]

\section{Software}\label{sec:software}
% §8.1 Binaries — 10 F2 binaries
% §8.2 Tests — 727 tests, 2 named regression locks
% §8.3 Provenance and empirical data — data/loop49/, W3C-PROV
[Software section goes here]

\section{Related work}\label{sec:related}
% §9.1 ML ablation methodology
% §9.2 Causal mediation
% §9.3 Sensitivity analysis (incl. Guo 2026)
% §9.4 Quantization (motivation for champion-scale follow-up)
[Related work goes here]

\section{Conclusion}\label{sec:conclusion}
% §10.1 Conclusion
% §10.2 Venue calibration (omit for submitted version; restore for
%        camera-ready or strip entirely as venue-specific meta)
[Conclusion goes here]

\subsubsection*{Acknowledgments}
% OMITTED FOR DOUBLE-BLIND REVIEW. Restore at camera-ready.

\bibliographystyle{tmlr}
\bibliography{f2_methodology}

\appendix

\section{Reproducible commands}\label{app:reproduce}
% §App. A.1-A.6
[Appendix A goes here]

\section{Provenance preamble format}\label{app:provenance}
% §App. B.1-B.3
[Appendix B goes here]

\section{Code-to-paper crosswalk}\label{app:crosswalk}
[Appendix C goes here]

\section{Test inventory}\label{app:tests}
% Pointer to papers/appendix_d_test_inventory.md (auto-generated)
[Appendix D goes here]

\end{document}
